\documentclass[11pt,a4paper]{report}

\usepackage[T1]{fontenc}
\usepackage[utf8]{inputenc}
\usepackage[margin=25mm]{geometry}
\usepackage{amsmath}
\usepackage{siunitx}
\usepackage{booktabs}
\usepackage{graphicx}
\usepackage[siunitx]{circuitikz}
\usepackage{enumitem}
\usepackage[hidelinks]{hyperref}

\sisetup{per-mode=symbol, detect-weight=true}

\title{GTO and IGCT Thyristors\\[2mm]
\large Device Principles, Switching Behaviour and Gate Driver Design}
\author{}
\date{}

\begin{document}

\maketitle
\tableofcontents

%======================================================================
\chapter{Introduction}
%======================================================================

The gate turn-off thyristor (GTO) and its direct descendant, the integrated
gate-commutated thyristor (IGCT), are the highest-power fully controllable
semiconductor switches available in a single die. They occupy the region of the
power-electronics landscape above the reach of a single IGBT module: blocking
voltages of \SIrange{4.5}{10}{\kilo\volt} and controllable turn-off currents of
several kiloamperes, at the cost of a switching frequency limited to a few
hundred hertz \cite{bernet2000,rodriguez2007multilevel}.

Both devices are four-layer, latching, bipolar structures. Their distinguishing
feature with respect to the conventional phase-control thyristor is the ability
to interrupt the anode current by extracting the stored charge through the gate
terminal. What separates the GTO from the IGCT is not the silicon alone but the
way in which that extraction is performed --- and therefore the gate driver.
The IGCT is, by definition, a gate-commutated thyristor (GCT) sold together with
its gate unit; the gate unit is part of the device, not an accessory
\cite{klaka1997igct,carroll1997igct}.

For this reason the gate driver is treated here as a full chapter
(Chapter~\ref{ch:gatedriver}) rather than as an application note. The design of
the gate unit determines the maximum controllable current, the failure mode, the
need for a turn-off snubber and, ultimately, the cost and volume of the whole
converter.

\section{Scope of the document}

Chapter~\ref{ch:devices} describes the structure and the physical operating
principle of the GTO and of the GCT/IGCT. Chapter~\ref{ch:characteristics}
covers static and dynamic characteristics, losses and safe operating area.
Chapter~\ref{ch:gatedriver} is dedicated to the gate driver: requirements,
topologies, sizing equations, layout, auxiliary power supply, monitoring and
protection, with a worked numerical example. Chapter~\ref{ch:applications}
summarises typical applications and compares the devices with high-voltage
IGBTs. Numerical values quoted throughout the text are representative orders of
magnitude; the data sheet of the specific device always takes precedence.

%======================================================================
\chapter{Device Structure and Operating Principle}
\label{ch:devices}
%======================================================================

\section{The four-layer structure}

Both devices are built on the classical $p^{+}npn^{+}$ thyristor structure,
usually described as two coupled bipolar transistors: a wide-base $pnp$
transistor (anode emitter, $n$-base, $p$-base) and a narrow-base $npn$
transistor (cathode emitter, $p$-base, $n$-base). The gate contacts the $p$-base.
Denoting by $\alpha_{pnp}$ and $\alpha_{npn}$ the common-base current gains, the
anode current is

\begin{equation}
I_{A} = \frac{\alpha_{npn} I_{G} + I_{leak}}{1 - (\alpha_{pnp} + \alpha_{npn})},
\end{equation}

so the structure latches when $\alpha_{pnp} + \alpha_{npn} \rightarrow 1$
\cite{baliga2008,lutz2011}. In the on state, high-level injection floods the
$n$-base with electron--hole plasma, which is the origin of the very low on-state
voltage drop of these devices and, at the same time, of their slow turn-off.

\section{The GTO thyristor}

\subsection{Cathode segmentation}

To turn the device off, the $npn$ transistor must be desaturated by removing
holes from the $p$-base through the gate. Since holes must travel laterally
under the cathode emitter, the cathode is not a single area but is divided into
several thousand narrow fingers (typically \SIrange{100}{300}{\micro\metre}
wide) interdigitated with the gate metallisation and connected in parallel. This
segmentation limits the lateral distance that carriers have to travel and makes
turn-off possible at all.

\subsection{Turn-off gain and current filamentation}

In a conventional GTO the negative gate current is significantly smaller than
the anode current. The turn-off gain

\begin{equation}
\beta_{off} = \frac{I_{TGQ}}{|I_{G,peak}|} \approx \frac{\alpha_{npn}}{\alpha_{pnp} + \alpha_{npn} - 1}
\end{equation}

is typically in the range \numrange{3}{5}. The device therefore turns off
\emph{in thyristor mode}: for a storage time of the order of
\SIrange{10}{25}{\micro\second} the structure remains regenerative while the
plasma is progressively squeezed. Because the individual cathode fingers are not
perfectly identical, the conducting area shrinks non-uniformly and the residual
current concentrates in a few fingers --- current filamentation --- which is the
dominant failure mechanism at turn-off. The maximum controllable turn-off
current $I_{TGQM}$ is set by this mechanism, not by thermal limits.

\subsection{Consequences: snubbers and minimum times}

Two practical consequences follow:

\begin{itemize}[nosep]
  \item A \emph{turn-off (dv/dt) snubber} is mandatory. A polarised RCD snubber
        with a capacitor of typically \SIrange{3}{6}{\micro\farad} for a
        \SI{4.5}{\kilo\volt}/\SI{3}{\kilo\ampere} device holds the anode voltage
        low while the current is still being commutated, keeping the device out
        of the filamentation region. The snubber is bulky, lossy and expensive:
        its dissipation, $\tfrac{1}{2}C V_{DC}^{2} f_{sw}$, often dominates the
        converter loss budget.
  \item \emph{Minimum on-time and minimum off-time} of the order of
        \SIrange{10}{100}{\micro\second} must be enforced by the control, because
        the plasma distribution must be allowed to settle before the opposite
        transition. This constrains the modulation strategy and effectively
        limits $f_{sw}$ to a few hundred hertz.
\end{itemize}

\subsection{Anode side variants}

Three anode-side arrangements are common: \emph{symmetric} (reverse blocking,
thick $n$-base, high on-state drop), \emph{asymmetric} (an $n$-buffer layer plus
anode shorts, thinner base, lower drop, but no reverse blocking capability, used
with an antiparallel diode in voltage-source converters), and \emph{reverse
conducting} (RC), where the freewheeling diode is integrated on the same wafer.

\section{The GCT and the IGCT}

\subsection{Hard drive: unity turn-off gain}

The GCT is a GTO whose gate is driven \emph{hard}: the gate unit sinks the entire
anode current, $\beta_{off} \approx 1$, in less than about one microsecond
\cite{gruening1996harddrive,klaka1997igct}. The whole cathode current is
commutated from the cathode into the gate \emph{before} the anode voltage starts
to rise. The $npn$ transistor is then completely cut off and the device
finishes turning off as an \emph{open-base $pnp$ transistor}: the structure is no
longer regenerative, so there is no plasma pinching and no filamentation.

The essential requirement is therefore not a large gate charge but a very high
gate-current slope. Since

\begin{equation}
\frac{\mathrm{d}i_{G}}{\mathrm{d}t} = \frac{V_{GK}}{L_{\sigma}},
\label{eq:didt}
\end{equation}

with the gate--cathode reverse voltage $V_{GK}$ limited to about
\SI{20}{\volt} by the avalanche breakdown of the gate--cathode junction, the
only remaining degree of freedom is the loop inductance $L_{\sigma}$. Achieving
\SIrange{6}{10}{\kilo\ampere\per\micro\second} requires $L_{\sigma}$ of a few
nanohenries, which is impossible with a cabled gate connection. Hence the
\emph{integration} of the gate unit with the device: the IGCT is a GCT plus a
coaxially bonded, multilayer-PCB gate unit.

\subsection{Silicon improvements}

The hard-drive concept is combined with several device-level improvements
\cite{steimer1999igct,lutz2011}:

\begin{itemize}[nosep]
  \item \emph{Transparent anode emitter}: a thin, lightly doped $p$ emitter whose
        injection efficiency depends on current density, replacing anode shorts;
        it reduces the tail current and improves the on-state/turn-off trade-off.
  \item \emph{Buffer layer} ($n$ field stop): allows a thinner $n$-base for a
        given blocking voltage, hence lower on-state voltage and lower stored
        charge.
  \item \emph{Free-floating silicon} in a press-pack housing: double-sided
        cooling, high thermal cycling capability, and the important property that
        the device fails as a stable short circuit (short-circuit failure mode,
        SCFM), which allows series connection with $n+1$ redundancy.
  \item \emph{Reverse-conducting} (RC-IGCT) and \emph{reverse-blocking}
        (RB-IGCT) variants, the latter used in current-source converters
        \cite{weber2000rbigct}.
\end{itemize}

\subsection{Snubberless operation}

Because turn-off is homogeneous, the RCD turn-off snubber can be removed. What
remains is a \emph{di/dt clamp}: a saturable or air-core inductor of a few
microhenries in the DC link branch, with its own freewheeling path, which limits
the turn-on current slope of the IGCT and the reverse recovery stress of the
opposing diode to about \SI{1}{\kilo\ampere\per\micro\second}. This is a single
component shared by a whole phase leg, in contrast with one snubber per switch.

\begin{figure}[htbp]
\centering
\begin{circuitikz}[scale=0.85, transform shape]
  % DC link source
  \draw (0,0) to[V, l=$V_{DC}$, invert] (0,4.5);
  \draw (0,4.5) -- (2,4.5);
  % di/dt limiting inductor in the positive rail
  \draw (2,4.5) to[L, l=$L_{cl}$] (5,4.5);
  % clamp network across the inductor
  \draw (2,4.5) -- (2,6) to[C, l=$C_{cl}$] (4,6) to[D*, l=$D_{cl}$] (5.5,6) -- (5.5,4.5);
  \draw (2,6) -- (2,7.2) to[R, l=$R_{cl}$] (4,7.2) -- (4,6);
  % main switch
  \draw (5,4.5) -- (5,4) to[thyristor, l_=IGCT] (5,2.5);
  % freewheeling diode
  \draw (5,2.5) -- (6.8,2.5) to[D*, l_=$D_{fw}$] (6.8,0) ;
  % load
  \draw (6.8,2.5) -- (8.8,2.5) to[L, l_=load] (8.8,0);
  % negative rail
  \draw (0,0) -- (8.8,0);
\end{circuitikz}
\caption{Simplified snubberless IGCT phase leg with $\mathrm{d}i/\mathrm{d}t$
clamp circuit ($L_{cl}$, $D_{cl}$, $C_{cl}$, $R_{cl}$). No individual RCD
turn-off snubber is required.}
\label{fig:clamp}
\end{figure}

%======================================================================
\chapter{Static and Dynamic Characteristics}
\label{ch:characteristics}
%======================================================================

\section{On-state and conduction losses}

In the on state the device behaves as a $p$-$i$-$n$ diode in series with a
junction. For loss calculation the usual piecewise-linear model is adequate:

\begin{equation}
v_{T}(i_{T}) = V_{T0} + r_{T}\, i_{T},
\qquad
P_{cond} = V_{T0}\, I_{T,avg} + r_{T}\, I_{T,rms}^{2}.
\end{equation}

Typical values for a \SI{4.5}{\kilo\volt} device at rated current are
$V_{T0}\approx\SIrange{1.2}{1.7}{\volt}$ and a total on-state drop of
\SIrange{2}{3}{\volt} at several kiloamperes --- markedly lower than an IGBT of
the same blocking class, and the main reason these devices survive in very high
power converters.

\section{Turn-on}

Turn-on is initiated by a positive gate pulse. Because the plasma spreads
laterally from the gate fingers at a finite velocity, the critical rate of rise
of on-state current $\mathrm{d}i/\mathrm{d}t$ must be limited externally; the
clamp inductor of Fig.~\ref{fig:clamp} performs this function. A hard, high
$\mathrm{d}i_{G}/\mathrm{d}t$ turn-on pulse improves the initial conducting area
and reduces turn-on losses and localised heating.

\section{Turn-off}

The turn-off waveform is conventionally divided into:

\begin{description}[nosep]
  \item[Storage time $t_{s}$] from the gate turn-off command to the beginning of
        the anode voltage rise. \SIrange{10}{25}{\micro\second} for a GTO, of the
        order of \SI{1}{\micro\second} for an IGCT, and much less dispersed in
        the latter case --- which is what makes series connection practical.
  \item[Fall time $t_{f}$] rapid collapse of the anode current, typically
        \SIrange{1}{2}{\micro\second}.
  \item[Tail current] slow recombination of the residual charge in the $n$-base,
        lasting several microseconds and contributing a substantial part of the
        turn-off energy, since it flows at full DC-link voltage.
\end{description}

The turn-off energy is obtained by integration over the whole transient:

\begin{equation}
E_{off} = \int_{t_{0}}^{t_{1}} v_{AK}(t)\, i_{A}(t)\, \mathrm{d}t .
\end{equation}

\section{Safe operating area and limits}

Three distinct limits must be respected: the maximum controllable turn-off
current $I_{TGQM}$ (filamentation limit, strongly gate-unit dependent), the RBSOA
in the voltage--current plane during turn-off, and the maximum junction
temperature. In addition, cosmic-ray induced single-event burnout imposes a
derating of the DC-link voltage with respect to the rated blocking voltage,
usually to \SIrange{50}{60}{\percent} of $V_{DRM}$ \cite{abb5sya2033}.

\section{Comparison}

\begin{table}[htbp]
\centering
\caption{Representative comparison of GTO, IGCT and HV-IGBT.}
\label{tab:compare}
\begin{tabular}{lccc}
\toprule
 & \textbf{GTO} & \textbf{IGCT} & \textbf{HV-IGBT} \\
\midrule
Blocking voltage        & \SIrange{4.5}{6}{\kilo\volt} & \SIrange{4.5}{10}{\kilo\volt} & \SIrange{3.3}{6.5}{\kilo\volt} \\
Turn-off current        & \SIrange{3}{6}{\kilo\ampere} & \SIrange{3}{6}{\kilo\ampere}  & \SIrange{0.6}{1.5}{\kilo\ampere} \\
Turn-off gain           & \numrange{3}{5}              & $\approx 1$                   & voltage driven \\
Storage time            & \SIrange{10}{25}{\micro\second} & $\approx\SI{1}{\micro\second}$ & --- \\
Turn-off snubber        & required                     & not required                  & not required \\
On-state drop           & low                          & low                           & higher \\
Typical $f_{sw}$        & \SIrange{200}{500}{\hertz}   & \SIrange{0.5}{1}{\kilo\hertz} & \SIrange{1}{2}{\kilo\hertz} \\
Gate drive power        & high                         & high                          & low \\
Failure mode            & short circuit (press-pack)   & short circuit (press-pack)    & module: open/rupture \\
\bottomrule
\end{tabular}
\end{table}

%======================================================================
\chapter{Gate Driver Design}
\label{ch:gatedriver}
%======================================================================

\section{Why the gate unit defines the device}

For a MOSFET or an IGBT the gate driver delivers a charge of the order of
microcoulombs and its main design problems are speed, isolation and noise
immunity. For a GTO or a GCT the gate unit is a \emph{power converter in its own
right}: it must sink kiloamperes for microseconds, with a current slope of
kiloamperes per microsecond, from a low-voltage rail, tens or hundreds of times
per second. The gate unit therefore determines:

\begin{itemize}[nosep]
  \item the maximum controllable turn-off current;
  \item whether a turn-off snubber is needed;
  \item the spread of the storage time, and hence the feasibility of series
        connection;
  \item a non-negligible share of the converter losses and volume.
\end{itemize}

This is why an IGCT is delivered as an inseparable assembly of GCT wafer,
press-pack housing and gate unit \cite{carroll1997igct,abb5sya2032}.

\section{Functional requirements}

A complete gate unit must provide the following functions.

\begin{enumerate}[nosep]
  \item \textbf{Turn-on pulse}: a positive gate current pulse with high
        $\mathrm{d}i_{G}/\mathrm{d}t$ and a peak of tens of amperes (GTO) to a few
        hundred amperes (GCT), lasting a few microseconds, to ignite as much
        cathode area as possible simultaneously.
  \item \textbf{Back-porch (on-state) current}: a continuous positive gate
        current maintained for the whole conduction interval, from a few amperes
        to a few tens of amperes depending on device size. It prevents partial
        de-latching at low anode current and improves $\mathrm{d}i/\mathrm{d}t$
        robustness.
  \item \textbf{Turn-off pulse}: a negative gate current whose amplitude equals
        the anode current (GCT) or a fraction of it (GTO), with the highest
        achievable slope.
  \item \textbf{Reverse bias hold}: a sustained negative gate--cathode voltage
        during the whole off interval, to keep the $npn$ emitter junction
        reverse biased and prevent spurious re-triggering by $\mathrm{d}v/\mathrm{d}t$.
  \item \textbf{Auxiliary supply}: isolated for the full converter voltage, and
        able to deliver the average gate power.
  \item \textbf{Command interface}: optical, with low and reproducible
        propagation delay.
  \item \textbf{Monitoring and protection}: supply undervoltage lockout, minimum
        on/off time enforcement, status feedback to the controller.
\end{enumerate}

\section{Sizing the turn-off stage}

\subsection{Loop inductance}

Equation~\eqref{eq:didt} is the central design equation. With $V_{GK}$ fixed by
the gate--cathode avalanche voltage (about \SIrange{20}{22}{\volt}), the required
slope maps directly onto a maximum admissible loop inductance:

\begin{equation}
L_{\sigma,max} = \frac{V_{GK}}{\left(\mathrm{d}i_{G}/\mathrm{d}t\right)_{req}} .
\end{equation}

For $V_{GK}=\SI{20}{\volt}$ and
$\mathrm{d}i_{G}/\mathrm{d}t=\SI{6}{\kilo\ampere\per\micro\second}$ this gives
$L_{\sigma,max}\approx\SI{3.3}{\nano\henry}$. Every element of the loop --- the
capacitor bank ESL, the MOSFET package inductance, the PCB planes, the contact to
the gate ring --- must be budgeted against this figure. In a GTO gate unit, where
$\beta_{off}\approx4$ is accepted, a loop inductance of a few hundred nanohenries
gives \SIrange{30}{40}{\ampere\per\micro\second}, which is the practical
difference between the two technologies.

\begin{figure}[htbp]
\centering
\begin{circuitikz}[scale=0.9, transform shape]
  \draw (0,0) to[V, l=$V_{GK}$, invert] (0,3);
  \draw (0,3) to[C, l=$C_{bank}$] (2.2,3);
  \draw (2.2,3) to[R, l=$R_{on}$, a=MOSFET bank] (5,3);
  \draw (5,3) to[L, l=$L_{\sigma}$] (7.2,3);
  \draw (7.2,3) -- (8.6,3) node[above right,font=\small]{gate} -- (8.6,2.2);
  \draw (8.6,2.2) to[thyristor, l=GCT] (8.6,-1.0);
  \draw (8.6,-1.0) node[below right,font=\small]{cathode} -- (0,-1.0) -- (0,0);
\end{circuitikz}
\caption{Equivalent circuit of the IGCT turn-off loop. The current slope is set
by $V_{GK}/L_{\sigma}$; the capacitor bank must supply the full anode current
with negligible voltage sag.}
\label{fig:loop}
\end{figure}

\subsection{Capacitor bank}

During turn-off the bank supplies the charge extracted from the device. For a
triangular gate-current pulse rising to $I_{TGQ}$ at slope $S$, the commutation
time is $t_{c}=I_{TGQ}/S$ and the charge is

\begin{equation}
Q = \tfrac{1}{2}\, I_{TGQ}\, t_{c} = \frac{I_{TGQ}^{2}}{2S},
\qquad
C_{bank} \ge \frac{Q}{\Delta V_{acc}} ,
\end{equation}

where $\Delta V_{acc}$ is the voltage sag that can be tolerated without
degrading the slope (typically \SIrange{5}{10}{\percent} of $V_{GK}$). The bank
is realised as a large parallel combination of low-ESL electrolytic and film or
ceramic capacitors distributed around the gate ring, so that the ESL of the bank
adds in parallel and the inductance budget is met.

\subsection{Switch bank}

The turn-off switch is a parallel array of low-voltage ($\approx\SI{30}{\volt}$),
low-$R_{DS(on)}$ MOSFETs --- typically several tens of devices --- arranged
symmetrically around the gate ring. The number $N$ follows from two constraints:

\begin{equation}
N \ge \frac{I_{TGQ}}{I_{D,pulse}},
\qquad
N \ge \frac{R_{DS(on)}\, I_{TGQ}}{\Delta V_{R}},
\end{equation}

the second ensuring that the resistive drop $\Delta V_{R}$ does not eat
significantly into the \SI{20}{\volt} available for the $\mathrm{d}i/\mathrm{d}t$.
Symmetry matters as much as the count: an asymmetric layout produces unequal
sharing and local overheating, and increases the effective loop inductance.

\subsection{Worked example}

Consider a \SI{4.5}{\kilo\volt} GCT with $I_{TGQ}=\SI{4}{\kilo\ampere}$,
$V_{GK}=\SI{20}{\volt}$ and a required
$S=\SI{6}{\kilo\ampere\per\micro\second}$.

\begin{itemize}[nosep]
  \item Loop inductance: $L_{\sigma}\le 20/6\times10^{9}=\SI{3.3}{\nano\henry}$.
  \item Commutation time: $t_{c}=4000/6\times10^{9}=\SI{0.67}{\micro\second}$.
  \item Extracted charge: $Q=\tfrac12\cdot4000\cdot0.67\times10^{-6}\approx\SI{1.3}{\milli\coulomb}$.
  \item Capacitor bank for $\Delta V_{acc}=\SI{2}{\volt}$:
        $C_{bank}\ge 1.3\times10^{-3}/2 \approx \SI{670}{\micro\farad}$.
  \item MOSFET count with $I_{D,pulse}=\SI{100}{\ampere}$ per device:
        $N\ge 40$; choosing $N=60$ with $R_{DS(on)}=\SI{4}{\milli\ohm}$ gives a
        bank resistance of \SI{67}{\micro\ohm} and $\Delta V_{R}=\SI{0.27}{\volt}$,
        which is negligible.
  \item Energy drawn from the bank per turn-off:
        $E \approx \tfrac12 L_{\sigma} I_{TGQ}^{2} + Q\,V_{GK}
        \approx \SI{26}{\milli\joule} + \SI{26}{\milli\joule}$, so at
        \SI{500}{\hertz} the turn-off contribution to the gate-unit power is of
        the order of a few tens of watts.
\end{itemize}

The result --- a few nanohenries, hundreds of microfarads, tens of MOSFETs ---
explains the physical form of an IGCT gate unit: an annular, multilayer printed
circuit board of \SIrange{15}{25}{\centi\metre} diameter, bolted concentrically
to the press-pack gate ring, with the return path immediately underneath the
outgoing path.

\section{The GTO gate unit}

Because the GTO operates with $\beta_{off}>1$, the switch and capacitor bank are
smaller, but three additional problems appear.

\begin{enumerate}[nosep]
  \item \textbf{Storage time control.} The long and dispersed storage time must
        be accounted for in the dead-time of the modulator, and cannot be
        compensated by the gate unit.
  \item \textbf{Coordination with the snubber.} The turn-off snubber must be
        correctly sized and low-inductance; the gate unit alone cannot guarantee
        the RBSOA.
  \item \textbf{Interlocks.} Minimum on-time and minimum off-time must be
        enforced in hardware inside the gate unit, not only in the modulator,
        because a violation during a fault sequence destroys the device.
\end{enumerate}

A classical GTO gate unit uses three separate paths sharing the gate terminal: a
positive pulse path (charged capacitor plus switch and series inductor to shape
the pulse), a back-porch path (a current-limited linear or switched source), and
the negative path (bank plus MOSFET array) with a hold-off supply for the reverse
bias.

\begin{figure}[htbp]
\centering
\begin{circuitikz}[scale=0.85, transform shape]
  % common cathode rail
  \draw (-4,-1.8) -- (7.5,-1.8);
  % turn-on branch
  \draw (-4,-1.8) to[V, l=$+V_{on}$, invert] (-4,5.0);
  \draw (-4,5.0) to[R, l=$R_{on}$] (-1.5,5.0) to[nos, l=$S_{on}$] (1.5,5.0);
  % back-porch branch
  \draw (-2.4,-1.8) to[V, l=$+V_{bp}$, invert] (-2.4,3.0);
  \draw (-2.4,3.0) to[R, l=$R_{bp}$] (-0.6,3.0) to[nos, l=$S_{bp}$] (1.5,3.0);
  % turn-off branch
  \draw (-0.9,-1.8) to[V, l=$-V_{GK}$] (-0.9,1.0);
  \draw (-0.9,1.0) to[C, l=$C_{bank}$] (0.4,1.0) to[nos, l=$S_{off}$] (1.5,1.0);
  % common node towards the gate
  \draw (1.5,5.0) -- (1.5,1.0);
  \draw (1.5,1.0) to[L, l=$L_{\sigma}$] (4.6,1.0);
  \draw (4.6,1.0) -- (6.0,1.0) node[above right,font=\small]{G} -- (6.0,0.4);
  \draw (6.0,0.4) to[thyristor, l=GTO] (6.0,-1.8);
\end{circuitikz}
\caption{Conceptual three-path GTO gate unit: turn-on pulse, back-porch source
and turn-off bank sharing the gate terminal.}
\label{fig:gtogate}
\end{figure}

\section{Auxiliary power supply}

The gate unit power consumption has three components:

\begin{equation}
P_{GU} = \underbrace{V_{bp} I_{bp} D}_{\text{back porch}}
       + \underbrace{(E_{on}^{G} + E_{off}^{G}) f_{sw}}_{\text{switching}}
       + \underbrace{P_{aux}}_{\text{logic, monitoring}} ,
\end{equation}

and typically ranges from a few tens of watts to a few hundred watts per device.
Two supply strategies are used:

\begin{itemize}[nosep]
  \item An isolated DC/DC converter fed from a low-voltage bus through a
        high-isolation transformer (partial discharge and creepage rated for the
        full converter voltage, plus the transient $\mathrm{d}v/\mathrm{d}t$
        common-mode stress of the switching leg).
  \item Self-supply from the device itself, harvesting energy from the anode
        voltage or from the clamp circuit during the off interval, with a
        start-up path.
\end{itemize}

In both cases the gate unit stores enough energy locally to perform at least one
safe turn-off after a supply failure, and signals the loss of supply to the
controller.

\section{Signal interface, isolation and synchronisation}

The command is transmitted over plastic or glass optical fibre. Two requirements
dominate:

\begin{itemize}[nosep]
  \item \textbf{Low and reproducible latency.} In series-connected stacks, the
        jitter between gate units directly translates into unequal voltage
        sharing during the transient. Fibre length matching and equal-latency
        receivers are standard practice.
  \item \textbf{Status feedback.} A return fibre reports supply status,
        gate-unit readiness and, in advanced units, an indication derived from
        the gate--cathode voltage that the device actually turned off. This
        feedback is what allows a controller to detect a failed device in a
        redundant series stack.
\end{itemize}

The common-mode transient immunity of the isolation barrier must be specified for
the $\mathrm{d}v/\mathrm{d}t$ seen at the device terminals, which for an IGCT
phase leg is of the order of \SIrange{1}{5}{\kilo\volt\per\micro\second}.

\section{Layout guidelines}

\begin{itemize}[nosep]
  \item Build the gate unit as an annular multilayer PCB coaxial with the
        press-pack; keep the go and return conductors on adjacent layers with
        minimum dielectric thickness to minimise $L_{\sigma}$.
  \item Distribute the capacitor bank and the MOSFETs uniformly around the ring;
        each MOSFET should see nominally the same path length.
  \item Keep the gate-unit ground reference tied to the device cathode, and treat
        every measurement referenced to it as floating at full converter
        potential.
  \item Provide thermal management for the MOSFET bank and the bank ESR: at
        several hundred hertz these components dissipate continuously.
  \item Route the optical receivers and the logic supply away from the pulse
        loop; the $\mathrm{d}i/\mathrm{d}t$ in the turn-off loop is the largest
        interference source in the whole converter.
\end{itemize}

\section{Protection and interlocks}

\begin{itemize}[nosep]
  \item \textbf{Undervoltage lockout} on the gate-unit rails, with inhibition of
        turn-on commands.
  \item \textbf{Minimum on-time / minimum off-time} enforced in hardware.
  \item \textbf{Turn-off inhibition window} during the storage phase, so that a
        command reversal cannot interrupt an ongoing transition.
  \item \textbf{Overcurrent handling.} A GCT cannot turn off an arbitrary fault
        current: above $I_{TGQM}$ the device will fail. Short-circuit protection
        is therefore performed at converter level (fault current limiting by the
        clamp inductor, crowbar or fuse coordination), not by the gate unit
        alone.
  \item \textbf{Failure reporting} through the return fibre, allowing the
        controller to trip the stack or, in redundant configurations, to continue
        operating with the failed device shorted.
\end{itemize}

%======================================================================
\chapter{Applications and Perspective}
\label{ch:applications}
%======================================================================

\section{Typical applications}

GTOs dominated high-power converters from the 1980s until the late 1990s, in
traction drives, rolling mill drives and early FACTS equipment. IGCTs replaced
them in most new designs from the late 1990s onwards
\cite{steimer1998hardgto,bernet2000}. Present application areas include:

\begin{itemize}[nosep]
  \item Medium-voltage industrial drives, typically three-level neutral-point
        clamped inverters in the \SIrange{1}{30}{\mega\watt} range
        \cite{rodriguez2007multilevel,wu2006};
  \item Marine propulsion and mine hoists;
  \item Grid interfaces: STATCOM, back-to-back interties, wind-farm converters;
  \item Pumped storage and large excitation systems;
  \item DC circuit breakers and fault current limiters, where the very high
        surge and turn-off current capability is decisive.
\end{itemize}

\section{Position with respect to IGBT and wide-bandgap devices}

The high-voltage IGBT has taken over a large part of the medium-voltage market
thanks to its simple voltage-controlled gate, higher switching frequency and
mature module technology. The IGCT retains an advantage where conduction losses
dominate, where the surge current capability and the short-circuit failure mode
are required, and in press-pack series stacks. Wide-bandgap devices --- SiC
MOSFETs and SiC IGBTs at \SIrange{3.3}{10}{\kilo\volt} --- are progressively
addressing the same voltage classes with much higher switching frequency, but
not yet with comparable single-die current ratings or surge capability.

\section{Conclusion}

The GTO and the IGCT differ far less in silicon than in the way the gate is
driven. Hard drive --- unity turn-off gain achieved through a gate loop of a few
nanohenries --- transforms a filamentation-limited, snubber-dependent switch into
a homogeneous, snubberless one, with a storage time reduced by an order of
magnitude and a dispersion low enough to allow series connection. Any design work
on these devices is therefore, to a large extent, gate-unit design work: loop
inductance budgeting, capacitor and switch bank sizing, auxiliary supply energy
management, and interlock logic.

%======================================================================
\bibliographystyle{unsrt}
\bibliography{gto_igct}

\end{document}
